% !TeX root = ../../../main.tex
\section{Conclusion}

% remember this is just the conclusion for the VisR transplant study, not the whole dissertation! no DoPIo here, and it's only for humans

This work investigated the feasibility of using VisR ultrasound elastography to noninvasively assess renal allograft fibrosis in human kidney transplant recipients. VisR-derived semiquantitative metrics of viscoelastic anisotropy (DoA) and heterogeneity (RR) were compared against clinically measured biomarkers and Banff scores from contemporaneous renal allograft biopsies. VisR metrics were found to mildly correlate with IFTA grades, especially in models using PD-based DoA and RR metrics involving the medulla. Classifiers using VisR metrics were able to differentiate RE- and RV-based DoA with mild IFTA as well as PD-based DoA and RR with moderate IFTA despite limited numbers of subjects with relevant biopsies for comparison. This study also identified some classifier models that predicted interstitial inflammation, inflammation in the presence of IFTA, and T-cell-mediated rejection, though classification performance was not significantly better than random guessing. Overall, these findings suggest that VisR has potential as a noninvasive tool for assessing renal allograft fibrosis, though further optimizations and larger clinical studies are necessary to validate its clinical utility.
